% =========================================================================
% Vol II Master Chapter: The seven faces of r(z) in 3d HT QFT
% =========================================================================

\chapter{The seven faces of $r(z)$ in 3d holomorphic-topological QFT}
\label{ch:dnp-identification-master}

The collision residue
$r(z)=\operatorname{Res}^{\mathrm{coll}}_{0,2}(\Theta_{\cA})$ of an
$\SCchtop$-algebra $\cA$ is the common genus-zero binary datum of
seven presentations: the bar--cobar twisting morphism, the DNP25
$A_\infty$ Yang--Baxter Maurer--Cartan element on line operators, the
classical $r$-matrix from the boundary PVA, the leading pole of the
GZ26 commuting Hamiltonians, the Yangian $r$-matrix, the Sklyanin
Poisson bracket, and the Gaudin Hamiltonian. On the common
comparison locus specified below the seven presentations realise the
same class. Genus, associator, level, Lie-type, and normalisation
hypotheses are theorem data.

\section{The 3d HT collision residue}
\label{sec:dnp-master-collision-residue}

Let $\cA$ be a logarithmic $\SCchtop$-algebra
\textup{(}Definition~\textup{\ref{def:log-SC-algebra}}\textup{)}. The bar complex
$\barB(\cA)$ on the configuration space $\FM_n(\bC)\times\Conf_n(\bR)$
carries a Maurer--Cartan element $\Theta_\cA = D_\cA - d_0$
\textup{(}Vol~I, Theorem~\textup{\ref{V1-thm:mc2-bar-intrinsic}}\textup{)}, whose
binary collision residue is the central object of this chapter:
\begin{equation}\label{eq:dnp-master-collision-residue}
r(z)
\;:=\;
\operatorname{Res}^{\mathrm{coll}}_{0,2}(\Theta_\cA)
\;\in\;
\bar{\cA}\,\widehat{\otimes}\,\bar{\cA}((z^{-1})).
\end{equation}
By the $\mathrm d\log$ absorption (the bar propagator
$\mathrm d\log(z-w)$ absorbs one power of $(z-w)$, so the
collision-residue pole order is one less than the OPE pole order),
the maximal pole order of $r(z)$ is $k_{\max}(\cA) = p_{\max}(\cA) - 1$, where
$p_{\max}$ is the maximal OPE pole among generating fields
\textup{(}Vol~I, Definition~\textup{\ref{V1-def:p-max}}\textup{)}. For the standard
landscape:
\begin{itemize}
\item Heisenberg: $p_{\max}=2$, so $r(z) = k/z$ (single pole).
\item Affine $\widehat{\fg}_k$: $p_{\max}=2$.  The trace-form
 presentation is $r_{\mathrm{tr}}(z)=k\,\Omega/z$.  The KZ/Sugawara
 presentation at non-critical level is
 $r_{\mathrm{KZ}}(z)=\Omega/((k+h^\vee)z)$.  The level-prefixed KZ
 representative is $r_{\mathrm{KZ,lev}}(z)=k\,\Omega/((k+h^\vee)z)$.
 These are $\mathrm{GRT}^{\mathrm{fin}}$-orbit representatives related
 by finite Casimir rescalings on the locus where the displayed scalars
 are defined (Remark~\ref{rem:grt-casimir-fm99}).  At $k=0$ the
 trace-form and level-prefixed representatives vanish; the
 KZ/Sugawara representative does not.  The relation is a tensor
 identity between representatives, not a scalar equality such as
 $k=1/(k+h^\vee)$.
\item $\beta\gamma$: $p_{\max}=1$, so $r(z) = 0$ in the classical sense
 (the simple pole is fully absorbed); the shadow depth $r_{\max}=4$
 arises from composite fields.
\item Virasoro: $p_{\max}=4$, so
 $r(z) = (c/2)/z^3 + 2T/z$ (cubic plus simple poles).
\item $\Walg_3$: $p_{\max}=6$, so $r(z)$ has poles up to $z^{-5}$ in the
 $W$-$W$ channel.
\end{itemize}

\section{Face 1: Swiss-cheese bar-cobar twisting}
\label{sec:face-1-bar-cobar}

The first face of $r(z)$ is the bar-cobar twisting morphism. The
chiral bar-cobar adjunction
(Theorem~\ref{thm:scchtop-bar-cobar-adjunction}) produces a
universal twisting morphism
$\pi_\cA \in \mathrm{Tw}(\barBch(\cA), \cA)$; the binary collision
residue is its degree-$2$ collision component:
\[
r(z)
\;=\;
\pi_\cA\big|_{\mathrm{degree}\,2,\,\mathrm{collision}}.
\]
This is the algebraic definition of $r(z)$ from which the other
faces descend.

\section{Face 2: Line operators as $\cA^!_{\mathrm{line}}$-modules}
\label{sec:face-2-line-operators}

The second face identifies line operators with modules over the
Koszul dual; the statement is
Theorem~\ref{thm:dnp-bar-cobar-identification}
(\S\ref{subsec:dnp-identification},
chapters/connections/line-operators.tex), in three parts:

\begin{enumerate}[label=\textup{(\roman*)}]
\item \textbf{Monoidal identification.} Under the equivalence
 $\cC_{\mathrm{line}} \xrightarrow{\sim}
 \cA^!_{\mathrm{line}}\text{-}\mathbf{mod}$, the meromorphic tensor
 product $\ell_1 \otimes_z \ell_2$ corresponds to the $R$-twisted
 tensor product $V_{\ell_1} \otimes_{R(z)} V_{\ell_2}$.
\item \textbf{Maurer--Cartan identification.} The DNP $A_\infty$
 Yang--Baxter MC element $r(z) \in \cA^! \otimes \cA^!$ is the
 collision residue $\operatorname{Res}^{\mathrm{coll}}_{0,2}(\Theta_\cA)$.
\item \textbf{Non-renormalization $=$ Koszulness.} The one-loop
 exactness of line-operator OPE \textup{(\cite[Thm.\,4.1]{DNP25})} is
 the $E_2$-collapse of the bar spectral sequence (chiral Koszulness).
\end{enumerate}

The line-operator category captures $r(z)$ from the 3d HT
boundary.

\begin{remark}[Three lanes: cohomology, strict chains, Lie type]
\label{rem:dnp-identification-lanes}
\index{Dimofte--Niu--Py!chain-level comparison}
\index{Dimofte--Niu--Py!type scope}
Theorem~\ref{thm:dnp-bar-cobar-identification} splits into three
lanes, each a hypothesis in
Theorem~\ref{thm:vol2-seven-faces-master}:
\begin{enumerate}[label=\textup{(\roman*)}]
\item \textbf{Cohomological and module lane.}
On the completed perturbative, boundary-linear, chirally Koszul
locus of Theorem~\ref{thm:lines_as_modules}, the line category is
equivalent to modules over $\cA^!_{\mathrm{line}}$, the DNP
meromorphic tensor product is the $R$-twisted tensor product, and
the DNP $A_\infty$ Yang--Baxter Maurer--Cartan element is the
binary collision residue. This is the cohomological content of
\cite[Thm.\,4.1, Thm.\,5.5, Thm.\,6.1]{DNP25} joined to
Theorem~\ref{thm:dnp-bar-cobar-identification} through the
$E_2$-collapse of the bar spectral sequence; no associator is
chosen.
\item \textbf{Strict chain lane.}
The strict identification of the meromorphic-tensor chain model
\cite{DNP25} with the $R$-twisted ordered-bar chain model fixes a
Drinfeld associator $\Phi\in\mathrm{GRT}_1(\bQ)$. The
Kontsevich--Tamarkin formality
$\varphi_\Phi : C_*(\FM_\bullet(\bC)) \xrightarrow{\sim} E_2$
transports the chain-level DNP model to the chain-level bar-cobar
model. Absent the choice, the canonical object is the
$\infty$-quasi-isomorphism class; strict chain-level equality lives
in the $\mathrm{GRT}_1(\bQ)$-torsor of
Remark~\ref{rem:grt-torsor-3d-ht}.
\item \textbf{Lie-type and evaluation lane.}
The perturbative DNP construction holds for a Lie algebra $\fg$
and matter representation $V$. The type restriction enters when
Face~2 is identified with the Drinfeld/Yangian/Gaudin faces. The
type-$A$ evaluation core
($\mathfrak{gl}_N$, $\mathfrak{sl}_N$;
Theorem~\ref{thm:recognition-completed-dg-shifted-yangian}) and
the affine non-critical KZ/Gaudin comparison cover the present
clause. Simply-laced and folded classical types require the
Dynkin-folding compatibility between line-operator OPE and
$R$-twisted coproduct; exceptional types require a
Guay--Regelskis--Wendlandt Yangian ambient \cite{GRW18} with a DNP
OPE comparison, outside this clause.
\end{enumerate}
Face~2 in Theorem~\ref{thm:vol2-seven-faces-master} is therefore
the cohomological/module lane on the completed perturbative
boundary-linear locus, or the strict chain lane after a fixed
associator, with the Drinfeld/Yangian/Gaudin comparison restricted
to the Lie-type lanes listed.
\end{remark}

\section{Face 3: Classical PVA on the 3d HT boundary at
genus~$0$}
\label{sec:face-3-classical-pva}

\begin{remark}[Genus-zero scope inherited from Vol~I]
\label{rem:face-3-genus-zero}
Face~3 carries a genus-$0$ clause. The Vol~I anchor
Theorem~\ref{V1-thm:kz-classical-quantum-bridge} establishes the
classical/quantum identification at tree level on $\bP^1$; the
Vol~II master Theorem~\ref{thm:vol2-seven-faces-master} states
the PVA comparison on $\cM_{0,n}$. Extension to genus $g\geq 1$
loads logarithmic Fulton--MacPherson clutching and the
KZB/elliptic $r$-matrix comparison as additional theorem data.
\end{remark}

The third face is the classical $r$-matrix from the boundary PVA
through the KZ25 3d holomorphic-topological Poisson sigma model.
The PVA $\lambda$-bracket on the boundary algebra $\cP$ produces
the classical $r$-matrix
\[
r^{\mathrm{cl}}(z)
\;=\;
\sum_{\alpha,\beta} c_0(e^\alpha,e^\beta)\,e_\alpha\otimes e_\beta\,/\,z
\;+\;\text{higher poles},
\]
via Gui--Li--Zeng quadratic duality \cite{GLZ22}. Vol~I,
Theorem~\ref{V1-thm:kz-classical-quantum-bridge}: at tree level on
$\bP^1$, $r(z) = r^{\mathrm{cl}}(z)$, and the gauge invariance of
the KZ25 sigma model is the $\lambda$-Jacobi identity, equivalent
to $d_{\barB}^2 = 0$ through the Arnold relation. Genus $g\geq 1$
demands the KZB connection, an elliptic or higher-genus
$r$-matrix, and clutching compatibility in logarithmic
Fulton--MacPherson compactifications; this face omits those
hypotheses.

For Heisenberg, affine KM, and Virasoro at genus~$0$, the
classical $r$-matrix from the PVA $\lambda$-bracket agrees with
the bar-complex collision residue at leading order; quantum
corrections enter class~M.

\section{Face 4: Bulk Hamiltonians from the sigma model}
\label{sec:face-4-bulk-hamiltonians}

The fourth face is the leading pole of the GZ26 commuting
differential operators on the genus-$0$ moduli space (Vol~I,
Theorem~\ref{V1-thm:gz26-commuting-differentials}). The shadow
connection
$\nabla^{\mathrm{hol}}_{0,n} = d - \operatorname{Sh}_{0,n}(\Theta_\cA)$
decomposes into commuting Hamiltonians
$\operatorname{Sh}_{0,n}(\Theta_\cA) = \sum_i H_i\,dz_i$; the
leading $1/z_{ij}$ coefficient in $H_i$ is $r(z_{ij})$ acting on
$V_j$. In 3d HT this is the dimensional reduction of the
sigma-model Hamiltonian to the genus-$0$ moduli space.

The operator order of $H_i$ is $k_{\max}-1$ for $k_{\max}\geq 1$
and $0$ for $k_{\max}=0$
(Theorem~\ref{V1-thm:shadow-depth-operator-order}):
\begin{itemize}
\item $k_{\max}=0$: trivial Hamiltonians ($\beta\gamma$).
\item $k_{\max}=1$: scalar/multiplication operators (Heisenberg, KM).
\item $k_{\max}\geq 3$: genuine differential operators of order $k_{\max}-1$ (Virasoro, $\Walg_N$).
\end{itemize}

\begin{remark}[Principal-$\cW_N$ normalisation obligation]
\label{rem:vol2-gz26-wn-normalization}
Existence and commutativity of the $\cW_N$ Hamiltonians from
$\Theta_{\cW_N}$ come from Vol~I,
Theorem~\ref{V1-thm:gz26-commuting-differentials}(v). Term-by-term
comparison with the operators of \cite{GZ26} requires an explicit
dictionary between principal $\cW_N$ generators, Calogero--Moser
Hamiltonian normalisation, and collision-depth coefficients of
$\Theta_{\cW_N}$. Face~4 carries existence, commutativity, and
leading-collision construction; term-by-term $\cW_N$ identification
awaits the normalisation dictionary.
\end{remark}

\section{Face 5: The dg-shifted Yangian}
\label{sec:face-5-dg-shifted-yangian}

The fifth face is the dg-shifted Yangian of \cite{DNP25}, with
$r(z)$ the universal $r$-matrix. The dg-shifted Yangian is the
boundary-colour Koszul dual $\cA^!_{\mathrm{line}}$ on the
chirally Koszul locus, its multiplication, coproduct, counit, and
$R$-matrix derived from the bar complex.

The Yangian quantisation parameter
\[
\hbar
\;=\;
\frac{1}{k+h^\vee}
\]
for affine $\widehat{\fg}_k$ (Drinfeld \cite{Drinfeld85}) equals
the KZ25 sigma-model coupling and the collision-residue prefactor:
the three-parameter identification of Vol~I,
Theorem~\ref{V1-thm:yangian-sklyanin-quantization}.

\section{Face 6: The Sklyanin Poisson bracket}
\label{sec:face-6-sklyanin}

The sixth face is the Sklyanin Poisson bracket on $(\fg^!)^*$
(Semenov-Tian-Shansky \cite{STS83}). Given a classical $r$-matrix
$r^{\mathrm{cl}}(z)$ satisfying CYBE, the Sklyanin bracket on
$C^\infty((\fg^!)^*)$ is
\[
\{f,g\}_{\mathrm{Skl}}(x)
\;=\;
\langle x,[df,dg]\rangle
\;+\;\langle r^{\mathrm{cl}},df\otimes dg\rangle.
\]
The Yangian $Y_\hbar(\fg)$ is the canonical deformation
quantisation of this Poisson structure at $\hbar = 1/(k+h^\vee)$:
the classical limit of Face~5.

\section{Face 7: The Gaudin Hamiltonians}
\label{sec:face-7-gaudin}

The seventh face is the Gaudin Hamiltonians of the integrable
system associated to the Yangian. By Vol~I,
Theorem~\ref{V1-thm:gaudin-yangian-identification}, for affine
$\widehat{\fg}_k$ the GZ26 commuting Hamiltonians are the Gaudin
Hamiltonians of $Y_\hbar(\fg)$ in the evaluation representation:
\[
H_i^{\mathrm{GZ}}
\;=\;
\frac{1}{k+h^\vee}\,
\sum_{j\neq i}
\frac{\Omega_{ij}}{z_i-z_j}
\;=\;
\frac{1}{k+h^\vee}\,H_i^{\mathrm{Gaudin}},
\]
the Feigin--Frenkel--Reshetikhin Gaudin model \cite{FFR94}. Higher
Gaudin Hamiltonians (collision residues at depth $k\geq 2$) extend
to the full $A_\infty$ Yangian structure.

\section{The seven-face master theorem in 3d HT}
\label{sec:dnp-master-seven-face}

\begin{theorem}[The seven-face master theorem in 3d HT]
\label{thm:seven-face-master-3d-ht}
\label{thm:vol2-seven-faces-master}
\index{seven-face master theorem!3d HT|textbf}
Let $\cA$ be a logarithmic $\SCchtop$-algebra in the standard 3d HT
landscape, and put
\[
r(z)=\operatorname{Res}^{\mathrm{coll}}_{0,2}(\Theta_\cA).
\]
Assume the following comparison data.
\begin{enumerate}[label=\textup{(H\arabic*)},start=0]
\item $\cA$ lies on the completed perturbative, boundary-linear,
chirally Koszul locus used in Theorems~\ref{thm:lines_as_modules} and
\ref{thm:dnp-bar-cobar-identification}.
\item The PVA, Hamiltonian, and Gaudin comparisons are made over
$\cM_{0,n}$.
\item For affine $\widehat{\fg}_k$ one has $k\neq -h^\vee$, so that
$\hbar=(k+h^\vee)^{-1}$ is defined and the KZ/Sugawara representative
is finite.
\item For strict chain models one fixed associator
$\Phi\in \GRT_1(\bQ)$ is used throughout.  At cohomology level no
associator is chosen.
\item The Drinfeld/Yangian/Gaudin comparison is taken on the
type-$A$ evaluation core, with simply-laced or folded classical types
included only after the line-OPE/folding compatibility of
Remark~\ref{rem:dnp-identification-lanes}.
\item Principal $\cW_N$ terms are interpreted through the
collision-depth construction; a term-by-term comparison with
\cite{GZ26} requires the normalization dictionary of
Remark~\ref{rem:vol2-gz26-wn-normalization}.
\end{enumerate}
Under these hypotheses the binary collision residue admits the
following seven presentations:
\begin{enumerate}[label=\textup{(F\arabic*)}]
\item The bar-cobar twisting morphism $\pi_\cA$
 \textup{(}Vol~I, Theorem~\textup{\ref{V1-thm:collision-residue-twisting}}\textup{)}.
\item The DNP $A_\infty$ Yang--Baxter Maurer--Cartan element on
 $\cA^!_{\mathrm{line}}$
 \textup{(}Theorem~\textup{\ref{thm:dnp-bar-cobar-identification}}\textup{)},
 in the cohomological/module lane, or in the strict chain lane after
 fixing $\Phi$.
\item The classical $r$-matrix from the boundary PVA $\lambda$-bracket
 via Gui--Li--Zeng quadratic duality
 \textup{(}Vol~I, Theorem~\textup{\ref{V1-thm:kz-classical-quantum-bridge}}\textup{)},
 over $\cM_{0,n}$.
\item The leading $1/z_{ij}$ coefficient of the GZ26 commuting
 Hamiltonians on $\cM_{0,n}$
 \textup{(}Vol~I, Theorem~\textup{\ref{V1-thm:gz26-commuting-differentials}}\textup{)};
 for principal $\cW_N$, this means the MC construction and
 commutativity, with term-by-term GZ26 matching only after the
 normalization dictionary in \textup{(}H5\textup{)}.
\item The Yangian $r$-matrix in the evaluation representation
 \textup{(}Drinfeld~\cite{Drinfeld85}\textup{)}; the identification
 holds away from $k = -h^\vee$ where the Yangian specialization
 parameter $\hbar = 1/(k + h^\vee)$ is well-defined.
\item The classical limit of the Yangian commutator equals the Sklyanin
 Poisson bracket on $(\fg^!)^*$
 \textup{(}Semenov-Tian-Shansky~\cite{STS83}\textup{)}.
\item The Gaudin Hamiltonian of the integrable system on $\cM_{0,n}$
 \textup{(}Feigin--Frenkel--Reshetikhin~\cite{FFR94};
 Vol~I, Theorem~\textup{\ref{V1-thm:gaudin-yangian-identification}}\textup{)}.
 The literal rescaling $H_i^{\mathrm{GZ}} = (1/(k+h^\vee))
 H_i^{\mathrm{Gaudin}}$ is singular at $k = -h^\vee$; the
 identification is stated away from the critical level.
\end{enumerate}
On the common locus specified by \textup{(H0)--(H5)}, these seven
presentations are equivalent presentations of the same binary collision
residue $r(z)$.
\end{theorem}

\begin{proof}
Face F1 is the binary collision residue under the bar--cobar
adjunction. Face F2 is
Theorem~\ref{thm:dnp-bar-cobar-identification} under the completed
perturbative boundary-linear hypotheses and the associator
convention of (H0), (H3). Face F3 is Vol~I,
Theorem~\ref{V1-thm:kz-classical-quantum-bridge} on $\cM_{0,n}$.
Face F4 is Vol~I,
Theorem~\ref{V1-thm:gz26-commuting-differentials}, with the
principal-$\cW_N$ normalisation separated in (H5). Faces F5--F7
are the Drinfeld, Semenov-Tian-Shansky, and
Feigin--Frenkel--Reshetikhin identifications at the non-critical
evaluation point $\hbar=(k+h^\vee)^{-1}$. Each construction maps
its face to $\operatorname{Res}^{\mathrm{coll}}_{0,2}(\Theta_\cA)$;
the seven presentations agree on the common locus.
\end{proof}

\begin{proposition}[Common-codomain commutativity of the seven-face diagram]
\label{prop:seven-face-sasaki-commutativity}
\index{seven-face master theorem!common-codomain commutativity|textbf}
\index{seven-face master theorem!diagram commutativity}
Let $\mathcal{A}$ be an $\SCchtop$-algebra in the standard 3d HT
landscape, restricted to the common locus $\cL$ of
Theorem~\ref{thm:vol2-seven-faces-master}. Write
$\phi_i: \mathrm{F}_i(\cA) \xrightarrow{\sim}
r(z)$ for the seven pairwise identifications
\textup{(}$i\in\{1,\dots,7\}$\textup{)} supplied by
Theorem~\ref{thm:vol2-seven-faces-master}. Then the complete graph
$K_7$ of face-to-face identifications
\[
\phi_{ij} := \phi_j^{-1}\circ\phi_i : \mathrm{F}_i(\cA)
\xrightarrow{\sim}\mathrm{F}_j(\cA),
\qquad i,j\in\{1,\dots,7\},
\]
commutes: for every triple $i,j,k\in\{1,\dots,7\}$,
$\phi_{jk}\circ\phi_{ij} = \phi_{ik}$ as isomorphisms of
$\fg\otimes\fg$-valued meromorphic one-forms on the configuration
space $\mathrm{Conf}_2(\bC)^*$.
\end{proposition}

\begin{proof}
The pairwise identifications $\phi_i$ of
Theorem~\ref{thm:vol2-seven-faces-master} share the codomain
$r(z) = \operatorname{Res}^{\mathrm{coll}}_{0,2}(\Theta_\cA)$
computed from the Maurer--Cartan element
$\Theta_\cA\in\mathrm{MC}(g^{\SCchtop}_\cA)$. The datum is unique
as a cohomology class in $H^2(g^{\SCchtop}_\cA)$ by Gui--Li--Zeng
quadratic duality \cite[Thm.~5.2]{GLZ24} once the conformal
structure on $\cA$ is fixed; its scalar normalisation on each
chart of $\cL$ is set by the chosen representative: the
affine-KM trace-form representative $k\Omega/z$ and the
level-prefixed KZ representative $k\Omega/((k+h^\vee)z)$ vanish at
$k=0$, the KZ/Sugawara representative $\Omega/((k+h^\vee)z)$ does
not. The finite $\GRT^{\mathrm{fin}}$ Casimir rescaling records
the change of representative. Every face $\mathrm{F}_i$ targets a
common object, and the star-centre factorisation
\[
\phi_{ij} = \phi_j^{-1}\circ\phi_i
\]
makes the $K_7$-diagram a cocycle-free groupoid on seven objects:
any connected diagram of isomorphisms into one codomain commutes
strictly. For any triangle $(i,j,k)$,
\[
\phi_{jk}\circ\phi_{ij}
\;=\;\phi_k^{-1}\circ\phi_j\circ\phi_j^{-1}\circ\phi_i
\;=\;\phi_k^{-1}\circ\phi_i
\;=\;\phi_{ik},
\]
on the common restriction locus $\cL$. At cohomology level the
associator action on the genus-zero binary datum factors through
the finite Casimir-rescaling quotient
$\mathrm{GRT}^{\mathrm{fin}}\cong\bQ^\times$
(Remark~\ref{rem:grt-torsor-3d-ht}); scalar rescalings commute.
For strict cochains the same associator $\Phi$ is fixed across the
diagram, so no new cocycle enters the triangle calculation. The
21-edge graph of face-to-face maps $\{\phi_{ij}\}_{i<j}$ is a
commuting groupoid on $\cL$ with bar-intrinsic face $\mathrm{F}_1$
as terminal vertex. The super-extension
(Remark~\ref{rem:grt-torsor-3d-ht}(iii)) replaces $\mathrm{GRT}_1$
by its super-analogue; the cocycle-freeness argument carries
through modulo Koszul signs.
\end{proof}

\begin{remark}[Seven representatives and four generators]
\label{rem:seven-vs-four-resolution}
The theorem fixes seven orbit representatives. Under the coarser
relation identifying the classical Drinfeld, Sklyanin, and Gaudin
presentations, four constructions generate the seven. Both counts
describe the same common-codomain groupoid.
\end{remark}

\begin{remark}[The $\mathrm{GRT}$-torsor of seven-face presentations]
\label{rem:grt-torsor-3d-ht}
\label{rem:grt-casimir-fm99}
The seven faces are seven orbit representatives under the free
transitive action of the Drinfeld Grothendieck--Teichm\"uller
group $\mathrm{GRT}_1(\bQ)$ on coordinatised presentations of the
universal collision residue (Vol~I standalone
\texttt{seven\_faces.tex}, Theorem GRT-face-torsor). Three
consequences for the Vol~II 3d HT statement:

\begin{enumerate}[label=\textup{(\roman*)}]
\item \textbf{Infinite face family.} Each associator
$\Phi\in\mathrm{GRT}_1(\bQ)$ produces a face $\mathrm{F}_\Phi$;
$\mathrm{F}_1$ is the torsor origin; the seven above are seven
choices of $\Phi$. Brown's motivic associator gives
$\mathrm{F}_8$ (MZV-valued; weight $\geq 3$ detects class-$M$
non-formality); Willwacher's graph-cocycle associator gives
$\mathrm{F}_9$ (operadic, graph-complex).

\item \textbf{Casimir rescaling as tensor identity.} On affine
Kac--Moody the $\mathrm{GRT}$-action on the genus-zero binary
datum factors through $\mathrm{GRT}^{\mathrm{fin}}\cong\bQ^\times$.
The trace-form $r_{\mathrm{tr}}(z)=k\,\Omega/z$, KZ/Sugawara
$r_{\mathrm{KZ}}(z)=\Omega/((k+h^\vee)z)$, and level-prefixed KZ
$r_{\mathrm{KZ,lev}}(z)=k\,\Omega/((k+h^\vee)z)$ representatives
are related by finite Casimir rescalings wherever the displayed
scalars are defined: a tensor identity between orbit
representatives, not the scalar equality $k = 1/(k+h^\vee)$. At
$k=0$ the trace-form and level-prefixed representatives vanish;
the KZ/Sugawara representative does not.

\item \textbf{$\mathrm{F}_7$ disambiguation.} Face F7 above is
$\mathrm{F}_7^{\mathrm{Gaudin}}$, the simple-pole projection
$r_1$, on all four shadow classes $G,L,C,M$. The class-$M$
top-$A_\infty$ face $\mathrm{F}_7^{M,\mathrm{top}}$ (the operation
$m_3$ detecting non-formality for Virasoro and $\Walg_N$) sits
outside the $\mathrm{GRT}$-orbit of $\mathrm{F}_1$: it vanishes on
classes $G,L,C$ and is written $\mathrm{F}_7'$ or
$\mathrm{F}_7^{M,\mathrm{top}}$ when ambiguity arises.

\item \textbf{Super-variant.} For chiral algebras with odd
generating fields, the seven-face theorem and
$\mathrm{GRT}$-torsor statement carry through after replacing
$\otimes$ by the Koszul-signed graded tensor product, $\Omega$ by
the super-Casimir, and $\mathrm{GRT}_1$ by its super-analogue on
the parity-graded pure-braid Lie superalgebra. The
$\mathrm{F}_1\Leftrightarrow\mathrm{F}_2$ identification factors
through the super-twist
$\sigma(a\otimes b) = (-1)^{|a||b|}b\otimes a$.
\end{enumerate}

On the affine-KM binary datum the motivic and operadic directions
project to the same finite Casimir-rescaling quotient; higher MZV
and graph-complex content enters only at higher collision depth or
class-$M$ non-formal data.
\end{remark}

\section{The Drinfeld--Kohno four-path comparison}
\label{sec:dnp-master-dk-four-path}

The Drinfeld--Kohno comparison relates KZ monodromy,
quantum-group Casimir eigenvalues, the Yangian evaluation
specialisation, and the Verlinde truncation on the integrable
$\widehat{\mathfrak{sl}}_2$ locus.

\begin{theorem}[The four-path Drinfeld--Kohno comparison]
\label{thm:dk-four-path}
\index{Drinfeld-Kohno!four-path comparison|textbf}
For affine $\widehat{\mathfrak{sl}}_2$ at integrable non-critical level
$k\in\bZ_{>0}$ and spins $0\leq j_i\leq k/2$, the following four
constructions give the same root-of-unity comparison datum: the first
three give the same rank-$1$ braiding eigenvalue ratio, and the
Verlinde construction gives the corresponding allowed fusion channels:
\begin{enumerate}[label=\textup{(\arabic*)}]
\item \textbf{KZ monodromy.} The monodromy of the Knizhnik--Zamolodchikov
 connection $\nabla^{\mathrm{KZ}} = d - (1/(k+2))\sum_{i<j}\Omega_{ij}\,d\log(z_i-z_j)$
 on $V_{j_1}\otimes\cdots\otimes V_{j_n}$ around the discriminant
 divisor. (This is the bar-coefficient realization; the connection
 $1$-form is $\sum r_{ij}\,dz_{ij}$.)
\item \textbf{Quantum Casimir eigenvalues.} The eigenvalues of
 $q^{C_2}$ on irreducible summands, where
 $q = \exp(\pi i/(k+2))$ and $C_2$ is the second Casimir.
\item \textbf{Yangian $R$-matrix at Drinfeld specialization.} The
 Yang $R$-matrix $R^Y(u) = u\cdot I + i P$ evaluated at the spectral
 parameter $u_D = \cot(\pi/(k+2))$, the \emph{Drinfeld
 specialization}.
\item \textbf{Verlinde truncation.} The Verlinde fusion ring at
 level $k$, with truncation at the Weyl chamber boundary determining
 which fusion channels survive.
\end{enumerate}
\end{theorem}

\begin{proof}
Paths (1) and (2) agree by the Drinfeld--Kohno theorem for the KZ
connection at $q=\exp(\pi i/(k+2))$. For path (3), Yang's rank-$1$
evaluation $R$-matrix has eigenvalue ratio
\[
\frac{u+i}{u-i}.
\]
At the Drinfeld specialisation $u_D=\cot(\pi/(k+2))$ this is
\[
\frac{\cot\theta+i}{\cot\theta-i}=e^{2i\theta},
\qquad \theta=\frac{\pi}{k+2},
\]
the $q^2$ phase entering the KZ braiding eigenvalue ratio. Path
(4) selects the surviving channels by the
$\widehat{\mathfrak{sl}}_2$ Verlinde formula: the same
root-of-unity truncation at the Weyl alcove boundary. Braiding
eigenvalue ratio and fusion support agree on the integrable locus.
\end{proof}

\begin{remark}[Reality of the Drinfeld specialisation]
\label{rem:dnp-master-drinfeld-real}
The Drinfeld specialisation $u_D = \cot(\pi/(k+2))$ is real. The
Yang $R$-matrix has eigenvalue ratio $(u+i)/(u-i)$; at
$u_D = \cot(\theta)$ the ratio is
$(\cot\theta + i)/(\cot\theta - i) = e^{2i\theta}$, the desired
phase $q^2 = e^{2\pi i/(k+2)}$. The braiding eigenvalues (Hecke
generators $q$ and $-q^{-1}$) differ from the KZ monodromy
eigenvalues ($q^{1/2}$ and $-q^{-3/2}$); the Drinfeld--Kohno
comparison matches monodromy with the quantum Casimir.
\end{remark}

\section{Fusion rules from the $R$-twisted coproduct}
\label{sec:dnp-master-fusion}

The Verlinde fusion rules for $\widehat{\mathfrak{sl}}_2$ at level
$k$ arise from the $R$-twisted coproduct on the evaluation modules
of the ordered bar complex. On the integrable locus
$0\leq \lambda,\mu,\nu\leq k$ the bar-coproduct fusion support
coincides with the Verlinde formula
\[
N^\lambda_{\mu\nu}
\;=\;
\sum_{\sigma}
\frac{S_{\mu\sigma} S_{\nu\sigma} \overline{S_{\lambda\sigma}}}{S_{0\sigma}}.
\]

\section{The Virasoro spectral $R$-matrix in closed form}
\label{sec:dnp-master-virasoro-r}

\begin{computation}[Virasoro spectral $R$-matrix on primary states]
\label{comp:dnp-master-virasoro-r-matrix}
\index{Virasoro!spectral $R$-matrix closed form}
For the Virasoro algebra at central charge $c$ acting on a primary state
of conformal weight $h$, the spectral $R$-matrix admits the closed form
\begin{equation}
\label{eq:dnp-master-virasoro-r-closed}
R(z)
\;=\;
z^{2h}\,\exp\!\left(-\frac{c}{4z^2}\right)
\;=\;
z^{2h}\sum_{k=0}^{\infty}\frac{(-c/4)^k}{k!\,z^{2k}}.
\end{equation}
After factoring out $z^{2h}$, the first four even coefficients are
$1$, $-c/4$, $c^2/32$, and $-c^3/384$.
Only even powers of $1/z$ appear (bosonic parity). The non-terminating
character is the genus-$0$ signature of class~M (infinite shadow depth).
\end{computation}

The Virasoro exponential structure is the 3d HT realisation of
the rank-$1$ Yangian $r$-matrix $R^Y(u) = u\cdot I + iP$, the
linear-in-$u$ structure of affine KM replaced by the exponential
because the underlying OPE has higher pole order.

\section{The $\Walg_3$ commuting Hamiltonians: rank-$2$ calculation}
\label{sec:dnp-master-w3-hamiltonians}

For $\Walg_3$ with $p_{\max} = 6$ and $k_{\max} = 5$, the GZ26
commuting Hamiltonians are differential operators of order
$4 = k_{\max} - 1$ on the conformal block space. The
collision-depth expansion reads
\[
H_i^{(\Walg_3)}
\;=\;
\sum_{j\neq i}
\sum_{k=1}^{5}
\frac{c_k^{(\Walg_3)}(V_j,\partial_{z_j})}{z_{ij}^k},
\]
with coefficient operators $c_k^{(\Walg_3)}$ determined by the
$\Walg_3$ structure constants. In the primary-state normalisation
the coefficient identities are:
\begin{itemize}
\item Depth-$1$ residue: $\partial_{z_j}/z_{ij}$
\item Depth-$2$ residue: $h_j/z_{ij}^2$ (conformal weight)
\item Depth-$3$ residue: $w_j/z_{ij}^3$ (spin-$3$ charge)
\item Depth-$4$ residue: $0$ (the depth-$4$ contribution vanishes)
\item Depth-$5$ residue: $(c/3)/z_{ij}^5$ (trivial on primaries)
\end{itemize}

The $\Lambda$ composite field, central to $\Walg_3$, acts on
primaries by
\[
\Lambda_0|h\rangle
\;=\;
(h^2 - 3h/5)|h\rangle,
\]
with roots $h = 0$ and $h = 3/5$, independent of $c$.

\section{The first rank-$2$ $\mathfrak{sl}_3$ Yangian $R$-matrix}
\label{sec:dnp-master-sl3-yangian}

\begin{computation}[The $\mathfrak{sl}_3$ Yangian $R$-matrix from the
ordered bar]
\label{comp:dnp-master-sl3-yangian}
\index{Yangian!$\mathfrak{sl}_3$ from ordered bar|textbf}
For affine $\widehat{\mathfrak{sl}}_3$ at level $k$, the Yang
$R$-matrix on the fundamental representation $V = \bC^3$ takes the form
\[
R^{\mathfrak{sl}_3}(u)
\;=\;
u\cdot I + P,
\]
with $P$ the permutation. The Casimir
$\Omega = (1/2)\sum_a \lambda_a\otimes\lambda_a$ in the Gell-Mann
basis acts on $V\otimes V = \mathbf{6}\oplus\mathbf{3}^*$ with
eigenvalues $1/3$ (symmetric) and $-2/3$ (antisymmetric). The
level-prefixed KZ representative on
$\Conf_2^{\mathrm{ord}}(\bC)$ is
\[
r^{\mathfrak{sl}_3}(z)
\;=\;
\frac{k\Omega}{(k+3) z},
\]
matching Yang's $R$-matrix on the fundamental evaluation
representation at rank $2$. The trace-form representative is
$k\Omega/z$; the KZ/Sugawara representative is $\Omega/((k+3)z)$.
\end{computation}

\begin{remark}[Fundamental--adjoint distinction]
\label{rem:dnp-master-fund-adj}
The Yang form $R(u) = u\cdot I + \Omega$ satisfies the spectral
Yang--Baxter equation on the fundamental representation, where
$\Omega = P - I/N$ for $\mathfrak{sl}_N$. On the adjoint, the
Casimir carries eigenvalues $\{2,0,-3,-6\}$ with multiplicities
$\{27,20,16,1\}$ matching
$\mathbf{8}\otimes\mathbf{8} = \mathbf{27}\oplus\mathbf{10}\oplus
\mathbf{10}^*\oplus\mathbf{8}\oplus\mathbf{8}\oplus\mathbf{1}$,
and the Yang form alone does not satisfy YBE; the universal
$R$-matrix of $Y(\mathfrak{sl}_3)$ enters. Rank $2$ is the first
rank where the distinction between Yang $R$ and the universal
Yangian $R$ surfaces.
\end{remark}

\section{Cross-volume bridge}
\label{sec:dnp-master-cross-volume}

The seven-face structure realises in three volumes:
\begin{itemize}
\item \textbf{Vol~I}: the holographic-datum master chapter for
 chiral algebras on curves. The proved theorems live there:
 gz26 commuting differentials, KZ classical-quantum bridge,
 Gaudin--Yangian identification, Yangian--Sklyanin
 quantisation, shadow-depth--operator-order. Vol~I faces
 $F_i^{\mathrm{V1}}$ map to the present
 $F_i^{\mathrm{V2}}$ through the 3d HT boundary reduction,
 with $F_i^{\mathrm{V1}}$ equal to the boundary-algebra face
 $F_i^{\mathrm{V2}}$ of the 3d HT slab carrying boundary
 algebra $\cA$.
\item \textbf{Vol~III}: the CY-holographic-datum master chapter
 for Calabi--Yau quantum vertex chiral groups; the
 Maulik--Okounkov $R$-matrix as face~$4$ for $K3\times E$, the
 affine super Yangian $Y(\widehat{\mathfrak{gl}}_1)$ as
 face~$5$ for toric CY3, the elliptic Sklyanin bracket as
 face~$6$ for toroidal CY.
\end{itemize}

The three chapters compose through the master concordance, the
trilingual dictionary between realisations. The genus-$1$
extension, the identification of the KZB connection with the
elliptic $r$-matrix and the elliptic Gaudin model for affine
KM, sits in the Volume~I genus-1 seven-faces chapter.

\section{Standing hypotheses: associator, Koszul dual, scope, slab}
\label{sec:dnp-master-hypotheses}

The seven-face master theorem in 3d HT carries four standing
conditions, each part of the theorem data; each pins the precise
content of Theorem~\ref{thm:vol2-seven-faces-master} and
Proposition~\ref{prop:seven-face-sasaki-commutativity} on the
common locus.

\begin{remark}[Cohomology versus cochain associator action]
\label{rem:dnp-master-associator-dichotomy}
Proposition~\ref{prop:seven-face-sasaki-commutativity} establishes
strict commutativity of the seven-face groupoid on the common
locus $\cL$, cocycle-freeness following from
$\GRT^{\mathrm{fin}}$ collapse of the $\GRT_1(\bQ)$-torsor on the
genus-zero binary datum. At the cochain level each face-to-face
transport $\phi_{ij}$ is associator-dependent: fixing
$\Phi\in \GRT_1(\bQ)$ at face $F_i$ produces a strict
representative $F_i^{\Phi}$, and the composition with $F_j^{\Phi}$
is strict because the same $\Phi$ is fixed throughout. Cohomological
face commutativity is associator-independent; cochain-level face
commutativity sits in the $\GRT_1(\bQ)$-torsor of
Remark~\ref{rem:grt-torsor-3d-ht}.
\end{remark}

\begin{remark}[$\cA^{!}_{\mathrm{line}}$ is an
$(\SCchtop)^{!}$-algebra]
\label{rem:dnp-master-koszul-dual-operad}
Face F2 identifies $\cC_{\mathrm{line}}$ with
$\cA^{!}_{\mathrm{line}}$-modules, where $\cA^{!}_{\mathrm{line}}$
is the boundary Koszul dual from the chiral bar-cobar adjunction.
The object $\cA^{!}_{\mathrm{line}}$ is an algebra over the
Koszul-dual operad
$(\SCchtop)^{!} \simeq (\Lie, \Ass, \text{shuffle-mixed})$,
distinct from $\SCchtop$. The $\SCchtop$-datum of the 3d HT slab
is the pair $(\cZ^{\mathrm{der}}_{\mathrm{ch}}(\cA), \cA)$, bulk
acting on boundary. Face F2 records the equivalence between
$\cC_{\mathrm{line}}$ and modules over an algebra in the
Koszul-dual operad.
\end{remark}

\begin{remark}[Concrete construction across the nine specialisation
fibres]
\label{rem:dnp-master-concrete-vs-abstract}
Theorem~\ref{thm:vol2-seven-faces-master} holds on the common
chirally Koszul comparison locus under the F2 type and
chain-level conditions above. The concrete construction of all
seven face-to-$r(z)$ identifications with explicit generators and
relations matching standard literature is complete for two
anchors: $\fsl_{2}$-Yangian (rational limit) and affine
Kac--Moody $V_k(\fg)$ at non-critical level. The other seven
fibres of the unified chiral quantum group are concrete on the
evaluation-generated core: the elliptic case carries the KZB
connection and classical $r$-matrix at the classical level,
full quantum-$R$ quasi-triangularity at finite $\hbar$ open; the
toroidal case sits at the formal-disk level
(Schiffmann--Vasserot CoHA, Miki). See
Remark~\ref{rem:concrete-vs-abstract-scope} of the
unified-chiral-quantum-group chapter.
\end{remark}

\begin{remark}[The slab as bimodule]
\label{rem:dnp-master-slab-bimodule}
The Dimofte $D|T|N$ slab, carrying all seven faces of $r(z)$ as
boundary and line-operator data, is a bimodule over a pair of
boundary chiral algebras
$(\cA_{\mathrm{L}}, \cA_{\mathrm{R}})$. The bulk
$\cZ^{\mathrm{der}}_{\mathrm{ch}}(\cA)$ acts on both boundaries;
the slab fibre functor computes the relative bar
$B(\cA_{\mathrm{L}}, \cZ^{\mathrm{der}}_{\mathrm{ch}}(\cA),
\cA_{\mathrm{R}})$ through the bimodule. The seven faces of
$r(z)$ inscribe on the diagonal
$\cA_{\mathrm{L}} = \cA_{\mathrm{R}} = \cA$, where the relative
bar collapses to
$B(\cA)\otimes_{\cZ^{\mathrm{der}}_{\mathrm{ch}}(\cA)} B(\cA)$
and the binary collision residue
$r(z) = \Res^{\coll}_{0,2}(\Theta_{\cA})$ reads the diagonal
bimodule $r$-matrix. Off-diagonal slabs with
$\cA_{\mathrm{L}} \neq \cA_{\mathrm{R}}$ carry an $r$-matrix in
$\cA_{\mathrm{L}}^{!} \otimes \cA_{\mathrm{R}}^{!}[\![z^{-1}]\!]$
mediating the categorified bimodule structure; this is the
ambient of the universal-holography master theorem.
\end{remark}

\begin{remark}[Three anchors at the climax]
\label{rem:dnp-master-climax-anchors}
Theorem~\ref{thm:vol2-seven-faces-master} sits at the
bar-intrinsic apex of three anchors:
\begin{itemize}
\item[(a)] \emph{Universal holography.} On the boundary-linear,
 physics-bridge prefactorisation locus, the seven-face master is
 the F1 face of the universal-holography functor
 $\Phi_{\mathrm{hol}}\colon \cA \mapsto T_{\cA}$
 (Theorem~\ref{thm:universal-holography-master}):
 $r(z) = \Res^{\coll}_{0,2}(\Theta_{\cA})$ is the boundary-to-bulk
 obstruction class on the diagonal slab bimodule of
 Remark~\ref{rem:dnp-master-slab-bimodule}. The five rung
 corollaries (Corollaries~\ref{cor:rung-heisenberg},
 \ref{cor:rung-affine-km}, \ref{cor:rung-virasoro},
 \ref{cor:rung-w-N}, \ref{cor:rung-w-infty}) carry the seven-face
 data through each rung of the $E_n$-topologisation ladder.
\item[(b)] \emph{Universal conductor.} The Volume~I conductor
 $K = c + c^{!}$ specialises to F1 as the canonical
 Casimir-rescaling-invariant scalar pinning the absolute
 normalisation of $r(z)$ on each $\GRT^{\mathrm{fin}}$-orbit. The
 conductor formula is the $\GRT^{\mathrm{fin}}$-invariant shadow
 of F1.
\item[(c)] \emph{Volume~III $\Phi$ functor.} The CY-to-chiral
 functor $\Phi$ produces seven-face data at every CY target: K3
 through the Mukai-orthogonal Yangian
 $Y_{\hbar}(\mathfrak{so}(4,20))$, with Hodge-parity refinement
 for $Y_{\hbar}(\mathfrak{so}(4|20))$ (scope $n=2$);
 $\bC^{3}$ through the affine Yangian
 $\Yhbar(\widehat{\fgl_{\infty}})$ as F5 of the
 Schiffmann--Vasserot CoHA $r$-matrix; toric CY${}_{3}$ through
 the Maulik--Okounkov $r$-matrix as F4 (scope $n=1$). The
 Volume~III master concordance carries the seven-face data at
 each $\Phi$-target.
\end{itemize}
\end{remark}
